% sections/01_introduction.tex -- Introduction and summary of findings.
% Written from Sections 2-6 and Appendix A; every finding points to the section that supports it and
% keeps that section's status tags and new/known markers. Novelty follows the citation audit
% notes/reviews/citation_check.md (part b) and the narrowed tags of App. A and Sec. 5: do not widen
% them. Two placeholders mark results still pending (saddle-dyn; regime-map round 2).
\section{Introduction and summary of findings}\label{sec:intro}

\subsection{Question and motivation}\label{sec:intro-question}

Tensor networks whose contraction pattern is a tree (\emph{tree tensor networks}, TTNs), among them
matrix product states (MPS, or tensor trains), hierarchical Tucker formats and binary trees, are used as
machine-learning models for classification and generative modelling
\citep{stoudenmire2016supervised,novikov2018exponential,han2018unsupervised,stoudenmire2018learning,%
cheng2019tree,liu2019machine,glasser2019expressive}, and their expressive power has been related to
that of deep networks \citep{cohen2016expressive,levine2018deep}. How they train is less understood.
Work on their trainability has focused on barren plateaus, which occur for global but not for local
losses of unitary-embedded MPS \citep{liu2022presence} and are absent from the energy optimization of
isometric MPS, TTNS and MERA \citep{miao2024isometric,barthel2025absence}. Which dynamics gradient
methods produce on a TTN, from which initializations, and why, has not been characterized
(\cref{sec:intro-prior}).

A TTN regression model, $f(x)=\ip{\W}{\feat1(x)\otimes\dots\otimes\feat\Ns(x)}$ with $\W$ contracted from
$\Kc$ cores (\cref{sec:setup-model}), is multilinear in the cores and homogeneous of degree $\Kc$ in
them, like a deep linear network. There, small initialization gives saddle-to-saddle (incremental)
learning, in which the components of the target are learned one at a time after plateaus that lengthen
as the initialization shrinks \citep{saxe2014exact,jacot2022saddletosaddle}, and large initialization
gives lazy training, in which the model stays close to its linearization
\citep{jacot2018neural,chizat2019lazy}. What does the tree add? Four things that a product of matrices
lacks. (i)~A component of the target can be \emph{local} on the tree: it can differ from what has been
learned on a few sites only, carried by a few cores. (ii)~The number of cores $\Kc$ and the depth of the
tree are different quantities (an MPS with $\Ns$ sites has depth $\Ns-1$, a binary tree $\log_2\Ns-1$).
(iii)~The bond dimension $\chi$ plays the role of width, but what it can carry is capped by the
dimension of the cut each bond makes. (iv)~Teachers can be genuinely tensorial, so that the correlation
$\ip{\W}{\Wt}$ has many local maxima on the sphere.

We ask, for gradient flow on the exact population loss of a teacher (\cref{sec:setup}): (Q1)~which
saddles are visited from small initialization; (Q2)~how the plateau durations scale with the
initialization scale $\isc$, the number of cores $\Kc$ and the bond dimension $\chi$; (Q3)~how the
position of a component on the tree sets its escape (\emph{locality}); (Q4)~which component is learned
first; (Q5)~which regimes occur as $\isc$, $\chi$, $\Kc$ and the balance of the cores vary, and where
their boundaries lie; (Q6)~what an input covariance changes; (Q7)~what the structure of the teacher
changes; and (Q8)~what discrete steps, gradient noise and finite samples change. Two further questions
concern practice: (B1)~the optimizers used for tensor networks (Adam, sweeps, Riemannian methods), and
(B5)~real data. The model is kept minimal (whitened Gaussian inputs unless stated, realizable teachers,
double precision; \cref{sec:setup-model} says why), so that theory can be compared with experiment
number by number. The main focus is the saddle-to-saddle regime (\cref{sec:s2s}).

\subsection{Summary of findings}\label{sec:intro-summary}

\paragraph{The answer in one paragraph.}
From small initialization, gradient flow on a TTN is saddle-to-saddle. It rests near critical points of
the problem with all bond ranks fixed, and each escape raises the bond ranks along a \emph{route}, a
connected set of cores that must move together \derived{} \empirical. The first escape takes a time
$\propto\isc^{-(\Kc-2)}$ \proven, whose prefactor one leading-order integration predicts for each
initialization \empirical. After it, the tree decides: the cheapest route has $k^*$ cores, the Steiner
subtree of the cores that carry the sites where the next component differs \proven, and the next plateau
is absent, logarithmic in $1/\isc$, or $\propto\isc^{-(k^*-2)}$ for $k^*=1$, $2$, $\ge3$ \derived{}
\empirical, so one teacher has different plateau laws on different trees \empirical; $k^*$ is the order
of the saddle under a spectral condition, and the order can be lower from $k^*=5$ on \proven. Which
component comes first is a race set by the initialization when two components differ on three or more
cores, and greedy as $\isc\to0$ when they differ on two \derived{} \empirical. Non-orthogonal teachers
make the saddle after the first escape strict and logarithmic on an MPS but leave it of higher order on
binary and hierarchical Tucker trees \derived{} \empirical. Raising $\isc$ moves training through a rich,
non-incremental window into the lazy regime \empirical, which exists only when every bond reaches its cut
dimension \derived{} \empirical. Discrete steps below a sharp stability limit keep the saddle sequence
\derived{} \empirical; gradient noise lengthens later plateaus and can trap \empirical; Adam removes the
$\isc$-dependence of the first escape \derived{} \empirical; ALS, DMRG and Riemannian gradient descent
remove the plateaus \derived{} \empirical. \Cref{tab:intro-summary} gives one row per question, and the
numbered list below states the main findings with their novelty.

\begin{table}[tbp]
  \centering
  \caption{Findings by research question. Status tags as in the conventions above; ``new'' means that
    no prior statement was found (\cref{sec:intro-prior}); binary4 is the binary tree with $\Ns=4$, a
    depth-3 linear network (\cref{rem:A-binary4-dln}). Details and scope in the sections named.}
  \label{tab:intro-summary}
  \footnotesize
  \setlength{\tabcolsep}{2.5pt}\renewcommand{\arraystretch}{1.04}
  \begin{tabularx}{\textwidth}{@{}>{\raggedright\arraybackslash}p{1.45cm}>{\raggedright\arraybackslash}X
      >{\raggedright\arraybackslash}p{1.95cm}>{\raggedright\arraybackslash}p{2.85cm}
      >{\raggedright\arraybackslash}p{1.05cm}@{}}
    \toprule
    question & finding & status & new or known & section \\
    \midrule
    Q1 saddle sequence & An odeco teacher is learned one component at a time; plateaus sit at the
      critical points with all bond ranks fixed (loss within $2\cdot10^{-2}$, relative); each event raises
      every bond rank by one (497 of 531 runs); complete sequences on four trees, $\Kc=3$--$7$ & \derived{} (prediction), \empirical & \knownresult[binary4] \newresult[MPS4, binary8,
      HT4] & \ref{sec:s2s-phenomenology} \\
    Q2 plateau scaling & First escape at $\isc^{-(\Kc-2)}T(Z)$ plus an offset, $T(Z)$ computed from each
      initialization before training; the norm of the initialization gives only a lower bound, low by $3$ to
      $3\cdot10^7$; weak dependence on $\chi$, none on $d$ at full bonds (binary4) & \proven{} (scaling, bound),
      \derived{} (offset), \empirical & \knownresult[exponent; the identity is the activation law of
      \citealp{kunin2025alternating}] \newresult[TTN form, tests] & \ref{sec:s2s-escape} \\
    Q3 locality & The next escape uses the smallest route with a nonzero local residual, $k^*$ cores: no
      plateau, $\ln(1/\isc)$ (strict saddle) or $\isc^{-(k^*-2)}$ (higher-order saddle); confirmed for $k\le5$
      on trees with $\Kc=3$--$15$; $k^*$ is the saddle order under a spectral condition, not always from $k^*=5$ on &
      \proven{} (route, order), \derived{} (times), \empirical & \knownresult[mechanism] \newresult[Steiner
      route, rates, order theorem, tree dependence] & \ref{sec:s2s-locality} \\
    Q4 learning order & Components differing on $m\ge3$ cores: a race set by the initialization,
      independent of $\isc$; on $m=2$ cores: greedy as $\isc\to0$, consistent with
      $P(\text{wrong})\propto\isc^{(1-\varrho)/\varrho}$ (non-adjacent cores); only then is a wrong order corrected, through a strict
      saddle & \proven{} (race), \derived, \empirical & \knownresult[binary4]
      \newresult[classification by $m$, race, rate] & \ref{sec:s2s-order} \\
    Q5 regime map & Saddle-to-saddle, rich, then lazy as $\isc$ grows (centred model); lazy only with full
      bonds; on binary4 both boundaries scale as $\chi^{-1/2}$, the lazy one as predicted ($\chi\ge32$); the
      saddle-to-saddle boundary rises with $\Kc$; a light root makes training lazy & \empirical,
      \derived{} (lazy constants), \conjecture{} (large $\Kc$) & \knownresult[in other models]
      \newresult[for TTNs] & \ref{sec:regime-map}, \ref{sec:other-lazy}--\ref{sec:other-large-chi} \\
    Q6 input covariance & A product covariance is a preconditioned whitened problem; aligned covariance
      reweights the race by $\Lambda_r\tw r$ and the landscape by $\Lambda_r^{1/2}\tw r$; misalignment and
      cross-site correlations make the MPS escape from the rank-one saddle logarithmic & \proven{} (reduction), \derived,
      \empirical & \newresult[for TTNs; two-layer analogue: \citealp{gidel2019implicit}] &
      \ref{sec:cov-aligned}--\ref{sec:cov-correlated} \\
    Q7 teacher structure & Plateaus at rank-constrained critical points (blends of components,
      generically not SVD truncations); escape from the rank-one saddle strict and logarithmic on MPS, of higher order on
      binary and HT trees; bottom-up growth on binary trees; star teachers forget a larger learned
      component & \proven{} (Hessians), \derived, \empirical & \newresult & \ref{sec:teacher-generic},
      \ref{sec:teacher-cp}--\ref{sec:teacher-star} \\
    Q8 steps, noise, samples & GD keeps the saddle sequence below $\lr_c=2/(\Kc\tw1^{2-2/\Kc})$, with
      event times shifted by a computable $O(\lr)$; SGD never shortened a plateau, lengthens later ones and
      can trap; small initialization recovers the teacher with its bond ranks from finite samples, the lazy
      limit cannot & \proven{} ($\lr_c$), \derived, \empirical, \conjecture{} (SGD law) &
      \knownresult[analogues] \newresult[TTN forms] & \ref{sec:other-optim} \\
    B1 optimizers & GD, heavy ball: as gradient flow; Adam, RMSProp: same order, first escape in a
      $\isc$-independent number of steps until $\epsilon$ restores $\isc^{-(\Kc-2)}$; ALS, DMRG,
      Riemannian GD: no plateaus & \derived, \proven{} (ALS, DMRG), \empirical & \newresult[for TTNs;
      Adam's $\epsilon$ law concurrently: \citealp{sabanayagam2026optimizer}] & \ref{sec:other-optimizers} \\
    B5 real data & MNIST: first route a sibling pair of a binary tree, gap $A/\isc+B$; later MPS routes
      mostly logarithmic; the pixel order sets the final loss, ranked before training by a route-amplitude
      share & \empirical{} (pre-registered test), \conjecture{} (mechanism) & \newresult &
      \ref{sec:real-data} \\
    \bottomrule
  \end{tabularx}
\end{table}
\paragraph{Main findings.}
Each item gives its evidence and its novelty: \emph{known} (with a reference), a \emph{TTN instance} of a
known mechanism (with the mechanism's reference), or \emph{new} (no prior statement found by the searches
of \cref{sec:intro-prior}).
\begin{enumerate}[label=\arabic*.,leftmargin=*,itemsep=3pt]
  \item \textbf{Saddle-to-saddle learning on genuine tensor networks} (\cref{sec:s2s-phenomenology}).
    From small fan-in initialization, gradient flow learns an odeco teacher one component at a time on
    binary4, MPS4, binary8 and HT4 (MPS8: the first two components). Each plateau is the critical point
    with all bond ranks fixed at the number of components learned, each event raises every bond rank by
    one, and the plateaus grow as $\isc^{-(\Kc-2)}$ (over $1.5$ decades of $\isc$ for $\Kc=3,4$, less than
    one for $\Kc=7,8$) at higher-order saddles, $\lambda_{\min}(\Hperp)\propto\isc^{\Kc-2}$ \empirical{}
    (predicted \derived). \emph{Novelty:} on binary4, a depth-3 linear network, this is known
    \knownresult[\citealp{saxe2014exact,jacot2022saddletosaddle,achour2024loss}]; on the genuine tensor
    networks it is new \newresult.
  \item \textbf{The first escape: a known exponent, a prefactor for each initialization}
    (\cref{sec:s2s-escape}). On the first plateau gradient flow is, to leading order, the flow of the
    correlation $g=\ip{\W}{\Wt}$: its direction converges to a balanced rank-one configuration while its
    norm blows up at the time $\isc^{-(\Kc-2)}T(Z)$, an integral along the spherical flow from the
    initialization direction $Z$ \proven, to which an $O(1)$ offset is added \derived. Computed before
    training, this predicts the first escape of every run (relative error $4\cdot10^{-8}$ to
    $2.2\cdot10^{-3}$ at the smallest $\isc$, 231 initialization--teacher pairs) and its first component
    (1,847 of 1,848 runs) \empirical. The time computed from the norm of the initialization alone is a
    lower bound \proven, low by factors from $3$ to $3\cdot10^7$ that grow with $d$, $\chi$ and $\Kc$
    \empirical. \emph{Novelty:} a TTN instance of known results: the exponent
    \knownresult[\citealp{jacot2022saddletosaddle,kumar2025early}], the blow-up identity, which is the
    activation law of alternating gradient flows \knownresult[\citealp{kunin2025alternating}], the
    norm-based time, which is the escape time from a maximizer \knownresult[\citealp{kumar2025towards}],
    and the balance of the escape direction \knownresult[linear chains: \citealp{kumar2025early},
    Lemma~5]. New \newresult[TTN form]: their evaluation on trees, the per-run tests, and the size of the
    gap to the norm-based time.
  \item \textbf{Later escapes follow the tree: the route rule} (\cref{sec:s2s-locality}). Once a
    component is learned, new parameters lower the loss only along routes, connected sets of cores whose
    new entries combine into a term of the residual at no Gauss--Newton cost. The cheapest route has
    $k^*$ cores; for two components that differ on the sites $D$ it is the Steiner subtree of the cores
    that carry $D$ \proven. The next plateau is then absent ($k^*=1$), logarithmic in $1/\isc$ with a head
    start at a strict saddle of rate $\tw2\tw1^{(\Kc-2)/\Kc}$ ($k^*=2$), or $\propto\isc^{-(k^*-2)}$ at a
    higher-order saddle ($k^*\ge3$) \derived. In about 11,600 runs on nine trees with $\Kc=3$--$15$ the
    exponent is $k-2$ for $k=3,4,5$ with named exceptions (for $k=6,7$ only on binary8 and HT4 with
    balanced initialization), the $k=2$ slopes and
    curvatures follow the predicted rate including its dependence on $\Kc$, and one teacher has a
    logarithmic plateau on MPS8, a $\isc^{-3}$ one on binary8 and one consistent with $\isc^{-5}$ on HT8
    \empirical. With three or four components the rule, applied at the saddle actually reached, predicts
    the plateau levels and their laws, also where a literal reading of the Steiner count fails
    \empirical{} (as an algorithm: \conjecture, \cref{rem:A-sequence}). \emph{Novelty:} a TTN instance of
    known mechanisms, dormant weights of lowest homogeneity degree
    \knownresult[\citealp{kumar2026learning,rawal2026theory}] and, in our model, descent curves through
    connected rank-deficient components, the smallest of which escapes first
    \knownresult[\citealp{furman2026benign}]; the $k=2$ head start is known in two-layer linear networks
    \knownresult[\citealp{zhang2026saddletosaddle}]. New \newresult: that the route is the Steiner subtree
    of the residual's sites, which can be far smaller than a rank-deficient component; its coupling and
    rate; the closed-form head start for $k^*\ge3$; and the measured dependence on the tree.
  \item \textbf{When the route gives the saddle order: a theorem} (\cref{app:saddle-order};
    \cref{sec:s2s-locality}). The route always descends at order $k^*$, so the saddle order is at most
    $k^*$. At a saddle with one learned component $j$ we prove a matching lower bound \proven{}
    \newresult: if the largest eigenvalue of a \emph{generalized piggyback operator}, which couples the
    parts of the residual living on sets of cores whose Steiner subtrees have fewer than $k^*/2$ cores,
    is below the learned weight $\tw j$, the order is exactly $k^*$ (\cref{thm:A-order-lower}). The
    condition always holds for $k^*=2$; for $k^*=3,4$ it reduces to the piggyback graph,
    $\lambda_{\max}(\Gamma)<\tw j$, above which the saddle is strict; and it holds for every two-component
    teacher learned in weight order, on every tree (\cref{cor:A-order}). From $k^*=5$ on it can fail,
    and then the order is lower (\cref{thm:A-order-sharp}): on MPS6 a teacher with $k^*=5$ and a positive
    semidefinite Hessian has order~4, through new entries that point at a non-residual pattern together
    with rotations of two learned cores; from that saddle gradient flow escapes in $\propto\isc^{-2}$
    instead of $\isc^{-3}$ and then forgets the learned component \derived{} (checked on this teacher from
    the saddle and on the seeds that visit it from small initialization). Scope: one learned component;
    saddles with two or more learned components are not covered. \emph{Novelty:} new; lower bounds on
    saddle orders were known for deep linear networks, order two or more
    \knownresult[\citealp{achour2024loss}, Thm~7], and for one parity saddle of a binary tree
    \knownresult[\citealp{furman2026benign}, Prop.~5.3].
    \placeholder{Pending (remote experiment saddle-dyn, running): one sentence on the pre-registered test
    of the theorem's dynamical prediction, escape times $\propto\varepsilon^{-(\text{order}-2)}$ from the
    saddle plus a perturbation of size $\varepsilon$ and the reduced exponent from small initialization,
    on the cases with order below $k^*$ and their controls.}
  \item \textbf{Strict or higher-order: the tree decides} (\cref{sec:s2s-locality,sec:teacher-generic}).
    With one learned component of a pattern teacher the saddle is strict iff a two-core route exists or
    the piggyback graph has $\lambda_{\max}(\Gamma)>\tw j$ (a wrong order, or a \emph{star} of smaller
    components sharing a core) \proven{} \newresult[for general trees; deep linear networks:
    \citealp{achour2024loss}]; otherwise $\lambda_{\min}(\Hperp)$ vanishes on the plateau like $\isc^{k-2}$
    \empirical. Non-orthogonal teachers (overlapping components, misaligned covariance, random TTN
    teachers, correlated sites) open two-core routes wherever two adjacent cores have free legs: on an MPS
    the escape from the rank-one saddle becomes strict and logarithmic, with slopes within $0.4\%$ (from
    the saddle) and $1.3\%$ (from random initialization) of the predicted rates, while binary and HT
    trees, whose internal cores carry no site, keep a higher-order saddle and an escape $\propto\isc^{-1}$
    \derived{} \empirical{} \newresult. The landscape bound of \citet{furman2026benign} allows order $\Kc$
    at such a point of an MPS; for generic teachers it is $2$.
  \item \textbf{Order of learning: a race or greedy} (\cref{sec:s2s-order}). As $\isc\to0$ the first
    component learned is the local maximum of $g$ on the sphere whose basin contains the initial direction
    \derived. When two components differ on sites carried by $m\ge3$ cores (\emph{tensor-like} pairs, such
    as every odeco pair on MPS, HT and binary8), each is such a maximum and the order is a race set by the
    initialization, independent of $\isc$: exact at leading order on MPS \proven, with
    $P(\text{larger first})=0.60$--$0.88$ at weight ratio $0.5$ over 23 settings, the same winner per seed
    at every $\isc$, more random as $\Kc$ grows and slowly greedier as $\chi$ grows \empirical. When they
    differ on $m=2$ cores (\emph{matrix-like}, e.g.\ on binary4), learning becomes greedy as $\isc\to0$,
    the wrong order dying out consistently with $\isc^{(1-\varrho)/\varrho}$ for non-adjacent cores
    \derived{} (heuristic) \empirical; only then is a wrong order corrected, through a strict saddle or by
    adding the larger component on a spare channel. A smaller component with a shorter route is learned
    first (\emph{locality beats weight}, \cref{sec:s2s-locality}). \emph{Novelty:} the greedy limit and
    the strict wrong-order saddle of binary4 are known
    \knownresult[\citealp{jacot2022saddletosaddle,achour2024loss}], and the same tensor-versus-matrix
    dichotomy is known in multi-spike tensor PCA \knownresult[\citealp{benarous2025stochastic}]; the
    classification by $m$, the race on trees, the order statistics and the power law in $\isc$ are new
    \newresult.
  \item \textbf{Failure modes: dead routes} (\cref{sec:s2s-failure,sec:other-stalled}). A sequence can
    end early because a route dies: under the conservation laws its end cores decay at a rate
    $\propto\isc^2$ while its escape needs a time $\propto\isc^{-(k-2)}$, so on routes of five or more
    cores the ends are decided long before the route could escape \derived{} (reduced model). Measured
    \empirical: no trap at $\Kc=3$; on MPS4 the outcome is fixed by the initialization direction; with the
    bond dimension equal to the teacher rank most runs on trees with $\Kc=7,8$ never learn the last
    component; with pattern teachers 398 of 5,682 fan-in runs stall, all with $k\ge4$. Balanced
    initialization removes the dead routes in the runs checked but not all stalls, spare bond channels
    help only where a route end is not capped, and the mechanism's criterion is strongly associated with
    death but neither necessary nor sufficient. \emph{Novelty:} new \newresult; spurious minima away from
    the minimum-norm set are known in the landscape of our model \knownresult[\citealp{furman2026benign},
    Thm~G.1], and exactly balanced initialization can fail in four-layer matrix factorization
    \knownresult[\citealp{luo2025global}].
  \item \textbf{The map of regimes} (\cref{sec:regime-map}; \cref{sec:other-lazy,sec:other-rich,%
    sec:other-stalled,sec:other-large-chi}). With the centred model, raising $\isc$ moves training from
    saddle-to-saddle through a rich, non-incremental window to lazy training \empirical. A lazy regime
    exists only when every bond reaches its cut dimension; below it the centred target is unrealizable
    and runs end on a floor \derived{} \empirical, the TTN form of a linearized model that fits only its
    tangent space \knownresult[\citealp{chizat2019lazy,ghorbani2019limitations}]. On binary4 both
    boundaries scale as $\chi^{-1/2}$; the lazy one is predicted, $\isc_{\rm lazy}\sqrt\chi=6.40$--$6.56$
    against $6.50$ \derived{} \newresult[the TTN constants], the saddle-to-saddle one is not derived. The
    saddle-to-saddle boundary rises with $\Kc$ ($0.17$ to $0.75$ for $\Kc=3$--$7$ at $d=2$, $\chi=4$),
    and stalled runs invade its side \empirical. A light root makes training lazy at an up to 100 times
    smaller function-space scale, a heavy one does not \empirical{} \derived{} (two-layer linear networks
    share the first effect, \knownresult[\citealp{kunin2024get,domine2025lazy}]). At large bond
    dimension the standard and NTK parameterizations become lazy, while a $\mu$P-type parameterization
    built on the injective norms of the cores keeps the kernel moving \derived{} \empirical.
    \emph{Novelty:} such maps are known for linear and diagonal networks
    \knownresult[\citealp{woodworth2020kernel,tu2024mixed,kunin2024get,domine2025lazy}]; for TTNs they
    are new \newresult, mapped in detail on one tree and with brackets only beyond seven cores.
    \placeholder{Pending (regime-map round 2, running): one sentence on the saddle-to-saddle boundary at
    large $\Kc$ from a realizable sweep at $d=2$ and $d=3$, and on whether $\isc_c^{\Kc}$ or
    $\isc_c^{\Kc}/\sqrt{d^{\Ns}}$ stays constant.}
  \item \textbf{Steps, noise and practical optimizers} (\cref{sec:other-optim,sec:other-optimizers}).
    Gradient descent can converge to a global minimum only if $\lr<\lr_c=2/(\Kc\tw1^{2-2/\Kc})$
    \proven; below it, it learns in gradient flow's order (765 of 765 runs) with event times shifted by a
    computable $\lr c_e$ \derived{} \empirical, and above it the learned component oscillates at the edge
    of stability. Gradient noise never shortened a plateau; it lengthens later ones and can trap the
    dynamics at a saddle where gradient flow is only slow \empirical{} (its law: \conjecture). Adam and
    RMSProp keep the order but escape the origin in a number of steps that does not depend on $\isc$,
    until $\epsilon$ restores the gradient-flow law \derived{} \empirical; ALS, DMRG and Riemannian
    gradient descent have no plateaus, and ALS and DMRG iterates do not depend on $\isc$ at all \proven.
    \emph{Novelty:} the stability limit is known for deep linear networks
    \knownresult[\citealp{mulayoff2020unique}]; the modified flow behind the event shift
    \knownresult[\citealp{barrett2021implicit}], the two-cycle at the edge of stability
    \knownresult[\citealp{ghosh2025learning,macdonald2025convergence}] and stochastic collapse
    \knownresult[\citealp{chen2023stochastic}] are known mechanisms; Adam's bounded first escape and its
    $\epsilon$ law were found concurrently in aligned deep matrix factorization
    \knownresult[\citealp{sabanayagam2026optimizer}]. New for TTNs \newresult: the stability limit on
    general trees, the event shift, the SGD delay and trap laws, the order kept by Adam, and the
    $\isc$-independence of sweeps.
  \item \textbf{Covariance, teacher structure and real data} (\cref{sec:covariance-teacher}). A product
    covariance turns the problem into a whitened one with a preconditioned flow \proven: with an aligned
    covariance the race is that of the whitened problem with weights $\Lambda_r\tw r$ and the landscape
    that with weights $\Lambda_r^{1/2}\tw r$, which can reverse the order \derived{} \empirical. Plateaus
    of random TTN teachers sit at critical points of the rank-constrained problem, not at SVD truncations
    where the two differ, and the bond ranks grow one bond at a time; binary students learn bottom-up;
    star teachers make the flow forget a larger learned component \empirical. On MNIST the first route
    after the escape on a binary tree is a pair of sibling cores with the $k=3$ gap law $A/\isc+B$, and a
    route-amplitude share computed from the data before training ranks pixel orders by their final loss
    in a pre-registered test, but not by their early speed \empirical; why the order matters is a
    \conjecture. \emph{Novelty:} the reweighting has a two-layer analogue
    \knownresult[\citealp{gidel2019implicit}]; the rest is new \newresult.
\end{enumerate}

\subsection{Relation to prior work}\label{sec:intro-prior}
% Sources: notes/literature.md v1.2 (sections 1.4, 2, 5) and notes/reviews/citation_check.md (part b,
% rows B1-B27). Pointers to theorems are given where the sections cite them.
The literature review behind this report (\path{notes/literature.md}, seven clusters) and an audit of
every citation and of possible prior work for every new result (\path{notes/reviews/citation_check.md};
about 270 searches in all, arXiv listings up to October 2026) found no paper that characterizes the
saddle-to-saddle dynamics of gradient flow on TTNs with dense cores and $\Kc\ge3$. The closest work, by
cluster, is the following.

\paragraph{Deep linear networks.} Exact solutions from decoupled (aligned) initialization
\citep{saxe2014exact,saxe2019mathematical}, incremental learning and its thresholds
\citep{gidel2019implicit,gissin2020implicit,arora2019implicit,li2021towards}, and saddle-to-saddle
dynamics from generic small initialization, where the first escape is proven and the later saddles are
conjectured \citep{jacot2022saddletosaddle}. The full sequence from small random initialization is
proven at depth two \citep{jiang2023algorithmic}, which covers the TTN with two cores; deeper
factorizations are solved for decoupled or identity-type initializations
\citep{saxe2014exact,chou2024gradient}. Critical points are classified as strict or non-strict
\citep{achour2024loss}, and balance is the conservation law of gradient flow
\citep{arora2018optimization,du2018algorithmic,kunin2021neural,marcotte2023abide}. The binary tree with
four sites is a depth-3 linear network (\cref{rem:A-binary4-dln}), so our binary4 results are checks.

\paragraph{Homogeneous models near the origin.} From small initialization the weights of a homogeneous
model first converge in direction to a KKT point of the correlation with the target and escape after a
time that grows as a negative power of the initialization \citep{kumar2025early,kumar2025towards}.
Dormant weights of the lowest homogeneity degree stay small for a time of the same order
\citep{kumar2026learning}; escapes that need $1$, $2$ or $r\ge3$ small layers take a time of order one,
logarithmic, or $\propto\epsilon^{-(r-2)}$ \citep{rawal2026theory}; alternating gradient flows describe
saddle-to-saddle dynamics by an activation law with dormant amplitudes carried across phases
\citep{kunin2025alternating}; weaker directions get a head start in two-layer linear networks
\citep{zhang2026saddletosaddle}; and the leap of \citet{abbe2023sgd} counts the coordinates that must be
learned at once. These results supply the mechanism of our first and later escapes; the tree decides
which cores are dormant, how many, and with which coupling.

\paragraph{The landscape of our model.} \Citet{furman2026benign} study exactly our model and loss, on any
tree, as a landscape: minimum-norm local minima are global; non-global critical points lack full Tucker
rank; at a minimum-norm critical point, every connected component of rank-deficient bonds gives a
descent curve whose order is its number of cores, and with several components ``the smallest core
delivers the least suppressed escape''; one parity saddle has a lower bound on its order; points off the
minimum-norm set can be spurious minima; and polynomial-size TTNs are hard for statistical-query
learners. They do not treat dynamics. Our routes can be proper subtrees of one such component, our lower
bound on the saddle order holds at every saddle with one learned component that meets a spectral
condition (\cref{app:saddle-order}), and our results are dynamical.

\paragraph{Tensor factorization.} Small initialization biases CP factorization towards low rank, learned
one rank-one term at a time \citep{razin2021implicit}, and hierarchical tensor factorization towards low
hierarchical rank \citep{razin2022implicit}; the latter shares the mode tree of our HT trees, but its
internal nodes are local CP factorizations, not dense cores (\cref{sec:setup-model}). Over-parameterized
orthogonal tensor decomposition learns the largest component first \citep{ge2021understanding}; lazy
training of a tensor decomposition needs many components \citep{wang2020beyond}; in deep CP
factorization the growth exponent of a component depends on the depth of its factors
\citep{hariz2022implicit}; the Tucker landscape has descents
of order 2, 3 and 4 at balanced points \citep{frandsen2022optimization}; and gradient descent on
tensor-train completion shows a long plateau followed by abrupt rank growth
\citep{milanesi2021implicit}. In multi-spike tensor PCA the recovery order is a greedy selection on
weight times initial correlation for tensors of order $p\ge3$ and follows the signal-to-noise ratios for
$p=2$ \citep{benarous2025stochastic,benarous2025permutation}, the dichotomy we find between tensor-like
and matrix-like pairs.

\paragraph{Lazy and rich training.} Lazy training and the tangent kernel
\citep{jacot2018neural,chizat2019lazy,lee2019wide}, its failure when the tangent space misses the target
\citep{ghorbani2019limitations}, the crossover read from the change of the kernel
\citep{geiger2020disentangling} and its coherent second-order part \citep{bordelon2022selfconsistent},
the kernel-to-rich transition of diagonal networks with an explicit saddle-to-saddle algorithm
\citep{woodworth2020kernel,pesme2023saddletosaddle,berthier2023incremental}, imbalance as a second knob
\citep{kunin2024get,domine2025lazy}, mixed and lazy-then-rich dynamics \citep{tu2024mixed,lyu2024dichotomy},
and the parameterizations of wide networks \citep{yang2021tensor,yang2024spectral}. For tensor networks,
Gaussian-process and tangent-kernel limits are known for MPS and tensor trains
\citep{guo2021neural,wesel2025tensor}, and factor-wise scaling rules for structured layers
\citep{qiu2024compute,potapczynski2024searching}; we found no map of regimes and no feature-learning
parameterization for TTNs.

\paragraph{Tensor networks for learning.} Beyond the models cited above, TTNs rest on the numerical
analysis of hierarchical formats, their manifolds and their alternating and sweeping algorithms
\citep{hackbusch2009new,grasedyck2010hierarchical,oseledets2011tensor,holtz2012manifolds,%
uschmajew2013geometry,holtz2012alternating,white1992density}, which also underlie Riemannian training
\citep{novikov2018exponential,steinlechner2016riemannian}. Their trainability has been studied through
barren plateaus \citep{liu2022presence,miao2024isometric,barthel2025absence}; the long plateaus here are
of another kind, since the gradient is small because the initialization is, not because it
concentrates over random initializations (\cref{sec:other-stalled}).

\paragraph{Steps, noise and optimizers.} The edge of stability and catapults
\citep{cohen2021gradient,lewkowycz2020large}; stability limits, two-cycles and the selection of flat
minima in deep linear and scalar networks
\citep{mulayoff2020unique,kreisler2023gradient,ghosh2025learning,macdonald2025convergence}; the modified
flow of gradient descent \citep{barrett2021implicit}; SGD as a stochastic differential equation, its
stability and stochastic collapse \citep{li2017stochastic,mulayoff2024exact,chen2023stochastic};
experiments with SGD in the saddle-to-saddle regime of deep linear networks
\citep{corlouer2026stochastic}; and, concurrently with this work, the plateau dynamics of Adam and other
optimizers in deep matrix factorization \citep{sabanayagam2026optimizer}.

\paragraph{What is new.} Known results that hold in our setting are marked \knownresult{} where they
appear and are shown in our setting: the exponent $\Kc-2$ of the first escape, its blow-up identity
and the balance of the escape direction (\cref{sec:s2s-escape}), the decoupled modes of aligned
initialization (\cref{sec:s2s-aligned}), the head start of two-core routes and the strict wrong-order
saddles of binary4 (\cref{sec:s2s-locality,sec:s2s-order}), the stability limit and edge-of-stability
cycle of gradient descent and, concurrently, Adam's $\epsilon$ law (\cref{sec:other-optim,%
sec:other-optimizers}). New, as far as these searches found, are: the route rule with its Steiner
subtree, rates and measured dependence on the tree; the theorem on the saddle order with its
counterexample; the classification of the learning order by the number of cores that separate two
components, and the race on trees; the contrast between MPS and binary trees for generic teachers;
the failure modes; the TTN constants and measurements of the regime map; and the tests on real data
(the novelty notes of the list above give the details). The searches are weakest for 2025--2026 work on the training of MPS
and tensor trains (\path{notes/literature.md}, section 1.2).

\subsection{Outline}\label{sec:intro-outline}

\Cref{sec:setup} fixes the model, the teachers, the initializations, the dynamics and the diagnostics
by which regimes are assigned. \Cref{sec:regime-map} maps the regimes over the initialization scale, the
bond dimension, the number of cores and the balance of the cores. \Cref{sec:s2s}, the main section,
treats the saddle-to-saddle regime: phenomenology, an exactly solvable case, the first escape, later
escapes and their locality (with the theorem on the saddle order), the order of learning, failure
modes, and how each parameter acts (\cref{tab:s2s-parameters}). \Cref{sec:other-regimes} treats the other
regimes (lazy, rich, stalled, large bond dimension) and what discrete steps, gradient noise, finite
samples and the optimizers used in practice change. \Cref{sec:covariance-teacher} adds input covariance,
teacher structure and real data, and \cref{sec:open} collects the open questions. \Cref{app:derivations}
derives the theory, \cref{app:experiments} gives the settings of every experiment, and
\cref{app:reproducibility} says how every figure is regenerated.
